\documentclass[a4paper,12pt]{article}

% ============================================================
% Кодировка, язык, шрифт
% ============================================================
\usepackage[T2A]{fontenc}
\usepackage[utf8]{inputenc}
\usepackage[russian]{babel}
\usepackage{helvet}
\renewcommand{\familydefault}{\sfdefault}
\usepackage{microtype}

% ============================================================
% Математика
% ============================================================
\usepackage{amsmath,amssymb,amsthm,mathtools}
\usepackage{icomma}

% ============================================================
% Страница и типографика
% ============================================================
\usepackage[
  a4paper,
  left=20mm,
  right=20mm,
  top=20mm,
  bottom=22mm
]{geometry}

\usepackage{setspace}
\onehalfspacing

\usepackage{parskip}
\setlength{\parindent}{0pt}

\usepackage{enumitem}
\setlist[itemize]{
  leftmargin=7mm,
  itemsep=2pt,
  topsep=3pt
}
\setlist[enumerate]{
  leftmargin=8mm,
  itemsep=5pt,
  topsep=4pt
}

% ============================================================
% Таблицы
% ============================================================
\usepackage{tabularx,booktabs}
\renewcommand{\arraystretch}{1.3}

% ============================================================
% Графика
% ============================================================
\usepackage{tikz}
\usetikzlibrary{calc}
\usepackage{tkz-euclide}

\tkzSetUpLine[line width=0.9pt,color=black]
\tkzSetUpPoint[color=accent,fill=accent,size=3]
\tkzSetUpLabel[color=black,font=\small]

% ============================================================
% Цвета
% ============================================================
\usepackage{xcolor}

\definecolor{accent}{HTML}{008080}
\definecolor{darkaccent}{HTML}{004D4D}
\definecolor{textgray}{HTML}{333333}
\definecolor{softgray}{HTML}{6B7280}
\definecolor{lightaccent}{HTML}{E8F3F3}

\color{textgray}

% ============================================================
% Ссылки
% ============================================================
\usepackage[
  colorlinks=true,
  linkcolor=darkaccent,
  urlcolor=darkaccent
]{hyperref}

% ============================================================
% Колонтитулы
% ============================================================
\usepackage{fancyhdr}
\setlength{\headheight}{15pt}

\pagestyle{fancy}
\fancyhf{}
\fancyhead[L]{\small\color{softgray}Стереометрия}
\fancyhead[R]{\small\color{softgray}Ксения Клыкова}
\fancyfoot[C]{\small\color{softgray}\thepage}
\renewcommand{\headrulewidth}{0.3pt}
\renewcommand{\headrule}{
  \hbox to\headwidth{\color{accent}\leaders\hrule height \headrulewidth\hfill}
}

% ============================================================
% Заголовки
% ============================================================
\usepackage{titlesec}

\titleformat{\section}
  {\Large\bfseries\color{darkaccent}}
  {\thesection.}{0.55em}{}

\titleformat{\subsection}
  {\large\bfseries\color{accent}}
  {\thesubsection.}{0.55em}{}

\titlespacing*{\section}{0pt}{14pt}{7pt}
\titlespacing*{\subsection}{0pt}{10pt}{5pt}

% ============================================================
% Служебные команды
% ============================================================
\newcommand{\plane}[1]{\ensuremath{#1}}
\newcommand{\belongs}{\ensuremath{\in}}
\newcommand{\notbelongs}{\ensuremath{\notin}}

\newcommand{\keyidea}[1]{%
  \par\medskip
  {\color{darkaccent}\bfseries Ключевая мысль.} #1
  \par\medskip
}

\newcommand{\warning}[1]{%
  \par\medskip
  {\color{darkaccent}\bfseries Важно.} #1
  \par\medskip
}

\newcommand{\method}[1]{%
  \par\smallskip
  {\bfseries\color{accent}#1}
  \par\smallskip
}

% ============================================================
% Документ
% ============================================================
\begin{document}

% ============================================================
% ТИТУЛЬНЫЙ ЛИСТ
% ============================================================
\begin{titlepage}
\thispagestyle{empty}

\begin{center}

\vspace*{25mm}

{\color{accent}\rule{0.22\linewidth}{1.1pt}}

\vspace{12mm}

{\Huge\bfseries Чек-лист}

\vspace{6mm}

{\LARGE\bfseries
Стереометрия:
третья аксиома\\[2mm]
и следствия из аксиом
}

\vspace{12mm}

{\large
пересечение плоскостей,
единственность плоскости,
математическая запись доказательств
}

\vspace{24mm}

{\large
Лёвин Артём Александрович\\[2mm]
\textbf{эксклюзивно для Ксении Клыковой}
}

\vfill

{\large \today}

\vspace{16mm}

\begin{tikzpicture}[x=1cm,y=1cm]
  \draw[
    color=accent,
    line width=1.2pt
  ]
  (0,0)
  .. controls (3.0,0.7) and (6.0,-0.7) ..
  (9.0,0)
  .. controls (11.0,0.45) and (13.2,0.25) ..
  (15.0,-0.1);
\end{tikzpicture}

\end{center}

\end{titlepage}

% ============================================================
% ОСНОВНОЙ ЧЕК-ЛИСТ
% ============================================================

\section{Что необходимо держать в памяти}

Сегодняшняя тема строится вокруг трёх базовых аксиом стереометрии.
Особенно важна третья аксиома, однако в доказательствах постоянно
используются также первая и вторая.

\subsection{Аксиома A1}

Через любые три точки, \textbf{не лежащие на одной прямой},
проходит плоскость, и притом только одна.

Если точки \(A,B,C\) неколлинеарны, то

\[
\exists!\,\alpha:
\qquad
A\in\alpha,\quad
B\in\alpha,\quad
C\in\alpha.
\]

Здесь символ

\[
\exists!
\]

читается как \textbf{«существует единственный»}.

\keyidea{
Если в задаче требуется доказать \textbf{существование и единственность
плоскости}, полезно искать три неколлинеарные точки.
}

\subsection{Аксиома A2}

Если две различные точки прямой принадлежат плоскости,
то вся прямая принадлежит этой плоскости.

Если

\[
A\in a,\qquad B\in a,\qquad
A\in\alpha,\qquad B\in\alpha,\qquad A\ne B,
\]

то

\[
a\subset\alpha.
\]

Иными словами:

\[
A,B\in\alpha
\quad\Longrightarrow\quad
AB\subset\alpha.
\]

\warning{
Для применения A2 нужны \textbf{две различные точки одной прямой},
принадлежащие плоскости.
}

% ------------------------------------------------------------

\section{Третья аксиома стереометрии}

\subsection{Формулировка}

Если две различные плоскости имеют общую точку,
то они пересекаются по прямой, проходящей через эту точку.

Все общие точки этих плоскостей лежат на этой прямой.

Пусть

\[
A\in\alpha,\qquad
A\in\beta,\qquad
\alpha\ne\beta.
\]

Тогда существует прямая \(a\), такая что

\[
\alpha\cap\beta=a,
\qquad
A\in a.
\]

То есть пересечение двух различных непараллельных плоскостей имеет вид

\[
\boxed{\alpha\cap\beta=a}.
\]

\subsection{Как читать запись}

Запись

\[
\alpha\cap\beta=a
\]

означает:

\begin{itemize}
  \item прямая \(a\) лежит в плоскости \(\alpha\);
  \item прямая \(a\) лежит в плоскости \(\beta\);
  \item каждая точка прямой \(a\) является общей для
        \(\alpha\) и \(\beta\);
  \item других общих точек вне прямой \(a\) у этих плоскостей нет.
\end{itemize}

\keyidea{
Если плоскости уже имеют хотя бы одну общую точку и известно,
что они различны, их общие точки образуют \textbf{целую прямую}.
}

% ------------------------------------------------------------

\subsection{Наглядная схема}

\begin{center}
\begin{tikzpicture}[scale=0.95]

  % Плоскость alpha
  \tkzDefPoint(-3.2,-1.4){P}
  \tkzDefPoint(2.4,-1.4){Q}
  \tkzDefPoint(3.3,1.2){R}
  \tkzDefPoint(-2.3,1.2){S}

  % Плоскость beta
  \tkzDefPoint(-1.0,-2.2){U}
  \tkzDefPoint(1.0,-2.2){V}
  \tkzDefPoint(1.0,2.2){W}
  \tkzDefPoint(-1.0,2.2){Z}

  \tkzDrawPolygon[color=softgray,line width=0.8pt](P,Q,R,S)
  \tkzDrawPolygon[color=accent,line width=0.9pt](U,V,W,Z)

  % Линия пересечения
  \tkzDefPoint(0,-1.8){M}
  \tkzDefPoint(0,1.8){N}
  \tkzDrawSegment[color=darkaccent,line width=1.6pt](M,N)

  % Точки на линии
  \tkzDefPoint(0,-0.55){A}
  \tkzDefPoint(0,0.75){B}
  \tkzDrawPoints(A,B)
  \tkzLabelPoints[right](A,B)

  \tkzLabelPoint[above left](S){$\alpha$}
  \tkzLabelPoint[above right](W){$\beta$}
  \tkzLabelPoint[right](N){$a$}

\end{tikzpicture}
\end{center}

На рисунке

\[
\alpha\cap\beta=a,
\qquad
A,B\in a.
\]

% ------------------------------------------------------------

\section{Как применять аксиому A3}

\subsection{Тип 1. Найти фигуру пересечения двух плоскостей}

Пусть различные плоскости \(\alpha\) и \(\beta\)
имеют общую точку \(A\).

По A3

\[
\alpha\cap\beta=a,
\qquad
A\in a.
\]

Следовательно, их пересечение является \textbf{прямой}.

\method{Шаблон ответа}

\[
\boxed{
\text{По A3 плоскости пересекаются по прямой,
проходящей через их общую точку.}
}
\]

% ------------------------------------------------------------

\subsection{Тип 2. Две общие точки двух плоскостей}

Пусть

\[
A,B\in\alpha,
\qquad
A,B\in\beta,
\qquad
\alpha\ne\beta.
\]

По A3 плоскости пересекаются по некоторой прямой \(a\):

\[
\alpha\cap\beta=a.
\]

Все общие точки двух плоскостей лежат на прямой их пересечения.
Поэтому

\[
A\in a,\qquad B\in a.
\]

Через две различные точки проходит единственная прямая, следовательно,

\[
a=AB.
\]

Итак,

\[
\boxed{\alpha\cap\beta=AB}.
\]

\keyidea{
Две общие точки позволяют не только установить факт пересечения
плоскостей, но и \textbf{однозначно определить прямую пересечения}.
}

% ------------------------------------------------------------

\subsection{Тип 3. Проверить принадлежность точки прямой пересечения}

Пусть

\[
\alpha\cap\beta=a,
\qquad
C\in\alpha,
\qquad
C\in\beta.
\]

Точка \(C\) является общей точкой плоскостей.
По A3 каждая общая точка двух плоскостей лежит
на их прямой пересечения.

Следовательно,

\[
\boxed{C\in a}.
\]

% ------------------------------------------------------------

\subsection{Тип 4. Доказать, что точка не принадлежит второй плоскости}

Пусть

\[
\alpha\cap\beta=a,
\qquad
M\in\alpha,
\qquad
M\notin a.
\]

Требуется доказать:

\[
M\notin\beta.
\]

Удобно применить доказательство от противного.

Предположим, что

\[
M\in\beta.
\]

Тогда \(M\) является общей точкой плоскостей
\(\alpha\) и \(\beta\).

По A3 каждая их общая точка принадлежит прямой пересечения \(a\).
Значит,

\[
M\in a.
\]

Но по условию

\[
M\notin a.
\]

Получено противоречие. Следовательно,

\[
\boxed{M\notin\beta}.
\]

\method{Схема доказательства от противного}

\[
\text{предположение}
\Longrightarrow
\text{следствие}
\Longrightarrow
\text{противоречие условию}
\Longrightarrow
\text{предположение неверно}.
\]

% ------------------------------------------------------------

\section{Важная тонкость: три плоскости}

Наличие общей точки у трёх плоскостей \textbf{не означает},
что все три плоскости обязательно пересекаются по одной общей прямой.

Например, возможна ситуация

\[
\alpha\cap\beta=a,
\qquad
\beta\cap\gamma=b,
\qquad
\gamma\cap\alpha=c,
\]

причём прямые \(a,b,c\) различны и проходят через одну общую точку \(O\):

\[
O\in a,\qquad O\in b,\qquad O\in c.
\]

\warning{
Одна общая точка трёх плоскостей не определяет общую прямую.
}

Если же у трёх плоскостей имеются \textbf{две различные}
общие точки \(A\) и \(B\), то прямая \(AB\) принадлежит
каждой из этих плоскостей. В этом случае все три плоскости
действительно содержат общую прямую \(AB\).

% ------------------------------------------------------------

\section{Следствие 1: прямая и точка вне неё}

\subsection{Формулировка}

Через прямую и точку, не лежащую на этой прямой,
проходит плоскость, и притом только одна.

Пусть

\[
M\notin a.
\]

Тогда

\[
\boxed{
\exists!\,\alpha:
\qquad
a\subset\alpha,
\qquad
M\in\alpha.
}
\]

% ------------------------------------------------------------

\subsection{Почему такая плоскость существует}

Выберем на прямой \(a\) две различные точки \(A\) и \(B\):

\[
A,B\in a,
\qquad
A\ne B.
\]

Так как

\[
M\notin a,
\]

точки \(A,B,M\) не лежат на одной прямой.

По A1 через три неколлинеарные точки проходит единственная плоскость:

\[
\exists!\,\alpha:
\qquad
A,B,M\in\alpha.
\]

По A2, поскольку

\[
A,B\in a
\quad\text{и}\quad
A,B\in\alpha,
\]

вся прямая \(a\) лежит в плоскости \(\alpha\):

\[
a\subset\alpha.
\]

Следовательно,

\[
a\subset\alpha,
\qquad
M\in\alpha.
\]

% ------------------------------------------------------------

\subsection{Почему плоскость единственная}

Предположим, что существуют две плоскости
\(\alpha\) и \(\beta\), каждая из которых содержит
прямую \(a\) и точку \(M\).

Тогда выбранные на \(a\) точки \(A\) и \(B\) вместе с \(M\)
принадлежат обеим плоскостям:

\[
A,B,M\in\alpha,
\qquad
A,B,M\in\beta.
\]

Но точки \(A,B,M\) неколлинеарны.

По A1 через них проходит \textbf{единственная} плоскость.
Следовательно,

\[
\alpha=\beta.
\]

Значит, требуемая плоскость действительно единственна.

\keyidea{
Существование и единственность здесь обеспечивает A1,
а принадлежность всей прямой плоскости --- A2.
}

% ------------------------------------------------------------

\subsection{Наглядная схема}

\begin{center}
\begin{tikzpicture}[scale=1.0]

  \tkzDefPoint(-3,-1.5){P}
  \tkzDefPoint(3,-1.5){Q}
  \tkzDefPoint(3.7,1.5){R}
  \tkzDefPoint(-2.3,1.5){S}
  \tkzDrawPolygon[color=softgray,line width=0.8pt](P,Q,R,S)

  \tkzDefPoint(-2,-0.5){A}
  \tkzDefPoint(2,-0.5){B}
  \tkzDrawLine[add=0.15 and 0.15,color=darkaccent,line width=1.2pt](A,B)

  \tkzDefPoint(0.7,0.7){M}
  \tkzDrawPoints(A,B,M)

  \tkzLabelPoints[below](A,B)
  \tkzLabelPoints[above](M)
  \tkzLabelPoint[above left](S){$\alpha$}
  \tkzLabelPoint[right](B){$a$}

\end{tikzpicture}
\end{center}

Здесь

\[
M\notin a,
\qquad
a\subset\alpha,
\qquad
M\in\alpha.
\]

% ------------------------------------------------------------

\section{Следствие 2: две пересекающиеся прямые}

\subsection{Формулировка}

Через две пересекающиеся прямые проходит плоскость,
и притом только одна.

Пусть

\[
a\cap b=\{M\}.
\]

Тогда

\[
\boxed{
\exists!\,\alpha:
\qquad
a\subset\alpha,
\qquad
b\subset\alpha.
}
\]

% ------------------------------------------------------------

\subsection{Доказательство}

Выберем

\[
A\in a,\qquad A\ne M,
\]

и

\[
B\in b,\qquad B\ne M.
\]

Так как прямые \(a\) и \(b\) различны и пересекаются только в точке \(M\),

\[
A\notin b,
\qquad
B\notin a.
\]

Следовательно, точки

\[
A,\ M,\ B
\]

не лежат на одной прямой.

По A1 существует единственная плоскость \(\alpha\), для которой

\[
A,M,B\in\alpha.
\]

Точки \(A\) и \(M\) принадлежат прямой \(a\), поэтому по A2

\[
a\subset\alpha.
\]

Точки \(B\) и \(M\) принадлежат прямой \(b\), поэтому по A2

\[
b\subset\alpha.
\]

Следовательно, обе прямые принадлежат одной плоскости.

Единственность следует из A1:
любая плоскость, содержащая \(a\) и \(b\),
содержит три неколлинеарные точки \(A,M,B\).

% ------------------------------------------------------------

\section{Следствие 3: две параллельные прямые}

Две различные параллельные прямые по определению лежат
в одной плоскости.

При этом плоскость, содержащая две различные параллельные прямые,
единственна.

Если

\[
a\parallel b,
\qquad
a\ne b,
\]

то

\[
\boxed{
\exists!\,\alpha:
\qquad
a\subset\alpha,
\qquad
b\subset\alpha.
}
\]

Для обоснования единственности можно выбрать две точки \(A,B\in a\)
и точку \(C\in b\). Так как прямые различны,

\[
C\notin a.
\]

Следовательно, \(A,B,C\) неколлинеарны,
и по A1 через них проходит только одна плоскость.

% ------------------------------------------------------------

\section{Как выбирать нужную аксиому}

\begin{table}[ht]
\centering
\caption{Ориентиры для доказательства}
\begin{tabularx}{\linewidth}{@{}p{0.25\linewidth}X@{}}
\toprule
\textbf{Что нужно получить} & \textbf{Что искать в условии} \\
\midrule

Единственную плоскость
&
Три неколлинеарные точки. Обычно применяется A1.
\\

Всю прямую в плоскости
&
Две различные точки этой прямой, принадлежащие плоскости.
Применяется A2.
\\

Прямую пересечения плоскостей
&
Две различные плоскости и хотя бы одну общую точку.
Применяется A3.
\\

Принадлежность общей точки прямой пересечения
&
Точка одновременно принадлежит обеим плоскостям.
Применяется A3.
\\

Непринадлежность точки второй плоскости
&
Удобно предположить обратное и применить A3.
\\

\bottomrule
\end{tabularx}
\end{table}

% ------------------------------------------------------------

\section{Математическая нотация}

Полезные символы:

\[
\begin{aligned}
A\in\alpha
&\quad &&\text{точка \(A\) принадлежит плоскости \(\alpha\)},\\
A\notin a
&\quad &&\text{точка \(A\) не принадлежит прямой \(a\)},\\
a\subset\alpha
&\quad &&\text{прямая \(a\) лежит в плоскости \(\alpha\)},\\
\alpha\cap\beta=a
&\quad &&\text{плоскости пересекаются по прямой \(a\)},\\
a\cap b=\{M\}
&\quad &&\text{прямые пересекаются в точке \(M\)},\\
\exists
&\quad &&\text{существует},\\
\exists!
&\quad &&\text{существует единственный}.
\end{aligned}
\]

\warning{
Греческими буквами \(\alpha,\beta,\gamma\) обычно обозначают
\textbf{плоскости}, а строчными латинскими \(a,b,c\) ---
\textbf{прямые}.
}

% ------------------------------------------------------------

\section{Типичные ошибки}

\begin{enumerate}
  \item
  \textbf{Забыть условие \(\alpha\ne\beta\).}

  Для совпадающих плоскостей утверждение A3
  о прямой их пересечения в таком виде неприменимо.

  \item
  \textbf{Считать, что одна общая точка трёх плоскостей
  даёт одну общую прямую.}

  Это неверно. Разные пары плоскостей могут пересекаться
  по разным прямым, проходящим через одну точку.

  \item
  \textbf{Путать принадлежность точки и прямой.}

  Корректно:

  \[
  A\in\alpha,\qquad a\subset\alpha.
  \]

  \item
  \textbf{Использовать A3 для пересечения двух прямых.}

  A3 относится к \textbf{двум плоскостям}.
  Для прямых используются свойства прямых и аксиомы A1--A2.

  \item
  \textbf{Ссылаться на A1, не проверив неколлинеарность.}

  Перед применением A1 необходимо обосновать,
  что три выбранные точки не лежат на одной прямой.

  \item
  \textbf{Писать только очевидный вывод без основания.}

  Даже если результат виден по рисунку, в доказательстве следует указать:

  \[
  \text{«по A1»},\qquad
  \text{«по A2»},\qquad
  \text{«по A3»}.
  \]

  \item
  \textbf{Считать рисунок доказательством.}

  Чертёж помогает увидеть идею.
  Само доказательство строится на аксиомах,
  определениях и логических выводах.
\end{enumerate}

% ------------------------------------------------------------

\section{Алгоритм работы с простой задачей на доказательство}

\begin{enumerate}
  \item Сделайте аккуратный чертёж.

  \item Выпишите, какие точки принадлежат каким прямым и плоскостям.

  \item Определите, что требуется доказать:
  \begin{itemize}
    \item существование плоскости;
    \item единственность;
    \item принадлежность прямой;
    \item принадлежность точки;
    \item линию пересечения.
  \end{itemize}

  \item Найдите подходящую аксиому.

  \item Запишите короткую цепочку:

  \[
  \text{условие}
  \Longrightarrow
  \text{аксиома}
  \Longrightarrow
  \text{вывод}.
  \]

  \item Если прямой путь неудобен, попробуйте
  \textbf{доказательство от противного}.
\end{enumerate}

\keyidea{
В начальной стереометрии сложность часто заключается не в том,
чтобы увидеть ответ, а в том, чтобы \textbf{строго обосновать очевидное}.
}

% ============================================================
% ТРЕНИРОВОЧНЫЙ БЛОК
% ============================================================

\clearpage

\section*{Тренировочный блок}
\addcontentsline{toc}{section}{Тренировочный блок}

\textbf{Путь А.}

Выполните задания письменно.
Для задач на доказательство обязательно указывайте,
какая аксиома или какое следствие используется.

\subsection*{Средний уровень}

\begin{enumerate}[label=\textbf{\arabic*.}]

  \item
  Различные плоскости \(\alpha\) и \(\beta\)
  имеют общую точку \(A\).
  Какой геометрической фигурой является
  \(\alpha\cap\beta\)?
  Что можно сказать о положении точки \(A\)?

  \item
  Известно, что
  \[
  A,B\in\alpha,\qquad A\ne B.
  \]
  Докажите, что
  \[
  AB\subset\alpha.
  \]

  \item
  Плоскости \(\alpha\) и \(\beta\)
  пересекаются по прямой \(a\).
  Известно, что
  \[
  C\in\alpha,\qquad C\in\beta.
  \]
  Докажите, что
  \[
  C\in a.
  \]

\end{enumerate}

\subsection*{Выше среднего}

\begin{enumerate}[label=\textbf{\arabic*.},start=4]

  \item
  Различные плоскости \(\alpha\) и \(\beta\)
  имеют две различные общие точки \(A\) и \(B\).
  Докажите, что
  \[
  \alpha\cap\beta=AB.
  \]

  \item
  Плоскости \(\alpha\) и \(\beta\)
  пересекаются по прямой \(a\).
  Точка \(M\) удовлетворяет условиям
  \[
  M\in\alpha,\qquad M\notin a.
  \]
  Докажите, что
  \[
  M\notin\beta.
  \]

  \item
  Дана прямая \(a\) и точка \(M\), причём
  \[
  M\notin a.
  \]
  Докажите, что существует единственная плоскость,
  содержащая одновременно прямую \(a\) и точку \(M\).

  \item
  Прямые \(a\) и \(b\) пересекаются в точке \(M\).
  Докажите, что через них проходит единственная плоскость.

  \item
  Различные плоскости \(\alpha\) и \(\beta\)
  имеют три общие точки \(A,B,C\).
  Докажите, что точки \(A,B,C\) лежат на одной прямой.

  \item
  Плоскости \(\alpha\) и \(\beta\) пересекаются по прямой \(a\).
  Точки \(P,Q\in\alpha\), причём
  \[
  P,Q\in\beta,\qquad P\ne Q.
  \]
  Докажите, что
  \[
  a=PQ.
  \]

\end{enumerate}

\subsection*{Сложный уровень}

\begin{enumerate}[label=\textbf{\arabic*.},start=10]

  \item
  Три различные плоскости
  \(\alpha,\beta,\gamma\)
  имеют две различные общие точки \(A\) и \(B\).

  Докажите, что существует прямая,
  принадлежащая всем трём плоскостям,
  и определите эту прямую.

  Затем объясните, почему наличие только одной общей точки
  трёх плоскостей было бы недостаточно для такого вывода.

\end{enumerate}

\end{document}